\section{Conclusion}
\label{sec:conclusion}

We have argued that the right target for AI assistance in empirical policy
research is not the unobservable causal effect but the \emph{identification
credibility} of the study that estimates it. ARGUS operationalizes this stance: a
bounded-agentic pipeline decomposes identification into auditable dimensions,
gathers and assesses in-paper evidence, and localizes risk in a transparent,
human-reviewable report, while confining model agency to two scoped, traceable
sub-tasks. We validate such auditors with synthetic flaw injection, which
supplies reproducible local ground truth without access to the true effect, and
we show that once the leakage making this evaluation circular is removed, a
keyword baseline is revealed to be insufficient---blind to flawed-but-plausible
evidence. That negative result is the point of departure for the doctoral
programme: reasoning over evidence adequacy, validated on a real annotated corpus
and with human experts in the loop. We believe auditing identification
credibility is both a tractable and a responsible role for AI in empirical
environmental-policy research.
